\documentclass{article}

\usepackage{tabularx}

\usepackage{booktabs}

\title{Problem Statement and Goals\\\progname}

\author{\authname}

\date{}

\input{../Comments.text}

\input{../Common.text}

\begin{document}

\maketitle

\begin{table}[hp]

\caption{Revision History} \label{TblRevisionHistory}

\begin{tabularx}{\textwidth}{llX}

\toprule

\textbf{Date} & \textbf{Developer(s)} & \textbf{Change}\\

\midrule

Date1 & Name(s) & Description of changes\\

Date2 & Name(s) & Description of changes\\

... & ... & ...\\

\bottomrule

\end{tabularx}

\end{table}

\section{Problem Statement}

\wss{You should check your problem statement with the

\href{https://smiths.github.io/capTemplate/Checklists/ProbState-Checklist.pdf}

{problem statement checklist}}.

\wss{You can change the section headings, as long as you include the required

information.}

\subsection{Problem}

Municipalities are responsible for ensuring that roads are maintained to an
acceptable standard during winter weather conditions. Winter maintenance
activities, such as snow clearing, plowing, salting, and other road treatments,
are performed across large municipal road networks to maintain safe and usable
transportation infrastructure. In many cases, these activities are performed by
contractors on behalf of municipalities. Once maintenance work has been
completed, municipalities need to determine whether the resulting road
conditions satisfy their required winter maintenance standards.

A major challenge is that records of maintenance activity do not necessarily
demonstrate the quality or completeness of the work that was performed. For
example, knowing that a winter maintenance vehicle travelled along a particular
road confirms that the vehicle was present at that location, but does not
confirm that the road was left in an acceptable condition. A road may still
contain excessive snow, ice, or other conditions that do not meet the required
level of service after a maintenance vehicle has completed its work. As a
result, municipalities need information about the actual condition of the road
following maintenance rather than relying only on records indicating that a
maintenance activity occurred.

Assessing completed winter maintenance work is also complicated by differences
in maintenance requirements between municipalities. Municipalities may define
different standards for what constitutes an acceptable post-maintenance road
condition. Consequently, the same observed road condition may satisfy the
requirements of one municipality while failing to satisfy those of another.
Determining whether winter maintenance work has been completed satisfactorily
therefore requires both an assessment of the resulting road condition and
consideration of the standards applicable to the location where the work was
performed.

The ability to verify the quality of completed work is important for both
municipal service management and contractor accountability. Municipalities need
reliable evidence to determine whether contracted winter maintenance services
have achieved the required level of service. Without a consistent method of
assessing post-maintenance road conditions, verification may require significant
manual review and may be difficult to perform consistently across large numbers
of road segments. This can make it difficult to identify inadequate maintenance
work and to determine whether contractors have fulfilled their obligations.

The verification of completed work is also related to contractor payment.
Municipalities need sufficient evidence that required maintenance standards have
been satisfied before determining whether completed work qualifies for payment.
A lack of clear and consistent information about post-maintenance road
conditions can make this decision more difficult and can create uncertainty for
both municipalities and contractors. Therefore, the central problem is to
provide an objective and consistent assessment of post-maintenance road
conditions that can be evaluated against location-specific municipal standards
and used to support verification of completed winter maintenance work.

\subsection{Inputs and Outputs}

The problem can be characterized using three primary high-level inputs. The
first input is information representing the observed condition of a road after
winter maintenance or treatment has been performed. This information provides
evidence of the state in which the road was left following the maintenance
activity. The second input is the geographic location associated with the
observed road segment. The location is necessary because the requirements used
to evaluate a road may vary between municipalities. The third input is the
applicable winter maintenance standard or level-of-service requirement
associated with that location.

The primary output is an assessment of the post-maintenance condition of the
road. This assessment identifies the observed road state in a form that can be
consistently interpreted. A further output indicates whether the assessed road
condition satisfies the winter maintenance standard applicable to its location.
Together, these outputs provide information that can be used to determine
whether the completed maintenance work meets the municipality's requirements.

The assessment also provides information that can support contractor payment
verification. If the resulting road condition satisfies the applicable
maintenance requirements, the completed work can be identified as satisfying
the conditions associated with payment. If the requirements are not satisfied,
the work can instead be identified for further review. The final payment
decision remains associated with the municipality's policies and contractual
requirements.

\subsection{Stakeholders}

The primary stakeholders are municipalities responsible for managing winter
road maintenance services. Municipal staff need to determine whether roads have
been maintained according to the standards established for their jurisdiction.
They require reliable and understandable information about post-maintenance road
conditions so that they can evaluate service quality, identify maintenance work
that may require additional attention, and support decisions related to
contractor performance and payment.

Winter road maintenance contractors are also major stakeholders. Contractors
are responsible for performing services such as plowing, clearing, salting, and
other winter road treatments according to contractual and municipal
requirements. Since assessments of completed work may affect whether their work
is considered satisfactory and whether it qualifies for payment, contractors
have an interest in assessments being consistent, traceable, and based on the
applicable municipal standards.

Municipal VU Consulting Inc. is a stakeholder and the supervising organization
for the project. The organization provides domain knowledge related to municipal
winter maintenance operations, service requirements, and the verification of
maintenance quality. Its involvement helps ensure that the problem being
addressed reflects practical municipal winter maintenance needs.

Municipal decision-makers and staff responsible for contract administration,
service monitoring, and winter maintenance operations are additional users of
the resulting assessment information. These stakeholders may use the
information to review maintenance performance and support operational and
contractual decisions.

\subsection{Environment}

The system is intended to operate within the context of municipal winter road
maintenance. The information being assessed originates from roads that have
undergone winter maintenance or treatment and is associated with specific
geographic locations. The operating context therefore includes municipal road
networks, winter maintenance vehicles, contractors, and municipal service
standards.

The physical environment can vary considerably. Road conditions may include
different amounts of snow, ice, slush, water, exposed pavement, and other
post-treatment conditions. Observations may also be collected during different
weather and lighting conditions, including snowfall, precipitation, overcast
conditions, daylight, darkness, and reduced visibility. Road appearance may
further vary because of differences in road surfaces, surrounding
infrastructure, traffic, and accumulated winter material. These variations form
part of the environment in which post-maintenance road conditions must be
assessed.

Municipal users are expected to review assessment information using standard
desktop or laptop computers. The user-facing software environment should
therefore be accessible using commonly available computing hardware and modern
web-browser software without requiring specialized computing equipment from
municipal staff.

The user environment is distinct from the development and processing
environment. Users should not be required to understand or directly interact
with the underlying technical components responsible for producing road-state
assessments. Instead, the user environment should present the resulting
information in a form suitable for reviewing road conditions, applicable
standards, maintenance compliance, and information relevant to contractor
payment verification.

\section{Goals}

\subsection{Recognition of Post-Treatment Road State}

Examine the state of a road after winter maintenance and/or treatment has been done on it,

and assign a label to the road state. This is useful for municipalities to assess the quality and completeness of work

that is being done with respect to winter maintenance.

\subsection{Assessment of Road State Based on Location}

Based on the location of the road, assign a label to the road state. Different municipalities have different standards for road maintenance, and the same road state may be considered acceptable in one municipality but not in another.

This goal is useful to ensure that results are consistent with pre-existing standards that municipalities have defined.

\subsection{Payment Verification Module}

Based on the state of the road and the location, determine whether a contractor should be paid for their work.

This goal is useful to ensure that the work that contractors do is being assessed in a timely manner.

\section{Stretch Goals}

\begin{itemize}

    \item \textbf{Contractor Payment Platform:} Build a platform that automatically sends payment to contractors once the AI verifies that a plowing or clearing pass met the required standard. This extends the core payment-verification module from calculating and flagging payments to actually distributing the payments.

    \item \textbf{Dispute and Appeal Workflow:} Provide a structured process for contractors to contest a flagged discrepancy or a rejected pass. Municipal staff can review the supporting imagery and per-segment records to uphold or overturn the decision, with every step logged for auditability.

    \item \textbf{Extended Verification for Sidewalks and Bike Lanes:} Extend the verification pipeline beyond roadways to cover sidewalks and bike lanes, adapting the surface-condition classification and required standards to each type of infrastructure.

    \item \textbf{Road Clearance Duration Estimation:} Estimate how long a road segment will stay clear after treatment, using weather data (such as temperature, precipitation, and wind) together with salting and treatment information, to help municipalities anticipate when re-treatment is needed.

\end{itemize}

\section{Custom Documents (CDs)}

\wss{For CAS 741: State whether the project is a research project. This

designation, with the approval (or request) of the instructor, can be modified

over the course of the term.}

\wss{For SE Capstone: List your custom documents (CDs).  Potential CDs include

usability testing, code walkthroughs, user documentation, formal proof,

GenderMag personas, Design Thinking, etc.  (The full list is on the course

outline and in Lecture 02.) Normally the number of CDs will be two (one per

term).  Approval of the CDs will be part of the discussion with the instructor

for approving the project. The CDs, with the approval (or request) of the

instructor, can be modified over the course of the term.}

\newpage{}

\section*{Appendix --- Reflection}

\wss{Not required for CAS 741}

\input{../Reflection.text}

\begin{enumerate}

    \item What went well while writing this deliverable?

    \item What pain points did you experience during this deliverable, and how

    did you resolve them?

    \item How did you and your team adjust the scope of your goals to ensure

    they are suitable for a Capstone project (not overly ambitious but also of

    appropriate complexity for a senior design project)?

\end{enumerate}

\end{document}